\section{Introduction}
\label{sec:intro}

Teams that ship retrieval-augmented or question-answering systems need to decide, for every generated answer, whether it can be trusted. The cheapest automatic check in common use is the cosine similarity between an embedding of the answer and an embedding of the evidence or a reference answer. It costs milliseconds, needs no prompt and is easy to threshold. It is also known to be the wrong instrument for the question being asked: a sentence and its negation, or a sentence with one number changed, are near neighbours in embedding space \citep{anschutz2023correct,cao2025semantic,steck2024cosine}, and recent work bounds what any embedding-based detector can certify about hallucination in retrieval-augmented generation \citep{sinha2025semantic}. Practitioners therefore escalate to natural-language inference (NLI) models or to an LLM acting as judge \citep{liu2023geval,zheng2023judging}, which are more expensive by orders of magnitude and are themselves imperfect \citep{godbole2025verify}.

The practical question is therefore not whether cosine is a sound measure of factuality---it is not---but \emph{when it is enough}: on which data, for which kinds of error, at which operating point, and whether a cheap stage can safely resolve part of the traffic so that a judge is called only where it is needed. Three gaps make this hard to answer from the literature. First, comparisons often pool evidence-grounded faithfulness with reference-based correctness, two targets that disagree whenever retrieval fails. Second, aggregate accuracy on natural benchmarks mixes many error types and answer lengths, so it does not tell an engineer which edits a given evaluator misses; perturbation studies exist \citep{gharsallah2026trust,nguyen2026accuracy} but typically compare the corrupted answer with the original, so a score drop can reflect \emph{any} rewording rather than the factual change. Third, thresholds are usually tuned and tested on the same distribution, while deployments move across domains, generators and answer lengths \citep{dubanowska2025representation}.

We study these questions with only open models on one CPU-only machine, and report measured rather than assumed costs. Our contributions are:

\begin{itemize}\setlength{\itemsep}{1pt}
\item \textbf{A target-separated comparison} of lexical overlap, two embedding models in three sentence-to-evidence configurations, BERTScore, a zero-shot NLI scorer, a feature-based risk estimator and two locally served LLM judges, on human-labelled faithfulness data (RAGTruth, FaithBench), human-labelled correctness data (EVOUNA-NQ) and a synthetic out-of-domain set (HaluEval-QA), with cluster-bootstrap intervals and thresholds fitted only on held-out calibration data.
\item \textbf{A matched-locus counterfactual test}: \NCfPairs\ single-fact flips of \NCfBases\ human-verified faithful answers in eight error types, with the flip and an LLM paraphrase placed in the \emph{same} sentence, so that the \emph{inversion rate}---how often an evaluator scores the factual error at least as high as a harmless rewording---isolates sensitivity to the fact from sensitivity to surface change.
\item \textbf{A threshold-transfer analysis} across five evaluation domains and within RAGTruth across generator families and answer lengths, testing whether a false-acceptance budget set on one domain holds on another.
\item \textbf{A cost-aware cascade} in which a cheap risk score accepts, rejects or escalates to a judge under calibration-time error budgets, compared against cosine gating and a random-escalation control, with latency measured on the same hardware for every component.
\end{itemize}

\noindent All hypotheses, primary metrics and tests were written down before any evaluator was run on test data; departures from that plan are listed in Appendix~\ref{app:deviations}. The LLM-judge experiments are a reduced-scale pilot on fixed, source-clustered subsets, and the counterfactual labels are rule-generated rather than human-annotated; we state both limits wherever the corresponding numbers appear.

Our main findings are these. (1)~On pooled RAGTruth, MiniLM cosine ranks answers worse than token precision against the evidence (AUC \PoolCosAuc\ vs.\ \PoolTokPrecAuc), and removing the cosine features from a feature-based risk model leaves its AUC unchanged. On short-answer correctness, checking whether the answer contains the reference beats cosine by a wide margin (balanced accuracy \EvContainBacc\ vs.\ \EvCosBacc). (2)~A single changed fact lowers cosine in \CfCosDrop\% of pairs, but with natural-data thresholds cosine turns an accepted answer into a rejected one in only \CfCosDetect\% of pairs. A factual error scores at least as high as a paraphrase of the same sentence in \CfCosInv\% of triplets for cosine and \CfNliInv\% for NLI. (3)~Every evaluator, including the learned risk model, broke a 5\% false-acceptance budget when its threshold was moved to several other domains, up to \TrRiskOffMax\%; most also exceeded it in one or two of their own domains. (4)~A risk-gated cascade beats judging everything with a 3B judge that rejects most faithful answers. On its pooled RAGTruth subset a 14B judge beats every cheap method used alone, but it is lenient on summaries. It is also the only evaluator whose decisions track single-fact edits (\CfJStrongDetect\% paired detection, small exploratory sample), so gating it trades judge time for accuracy.
